% Part: normal-modal-logic
% Chapter: axioms-systems
% Section: derived-rules

\documentclass[../../../include/open-logic-section]{subfiles}

\begin{document}

\olfileid{nml}{prf}{der}

\olsection{Aturan Turunan}

Nggoleki lan nulis !!{derivation}s pancen angel,
mbulet, lan bola-bali. Upamane, kerep banget kita pengin ngalih saka
$!A \lif !B$ menyang $\Box !A \lif \Box !B$, yaiku ngetrapake
aturan~\RK. Iki mbutuhake panganggone \Nec, banjur nyathet instansi
saka~\Ax{K} kang cocog, banjur ngetrapake~\MP. Kanggo ngalih saka
$!A \lif !B$ lan $!B \lif !C$ menyang $!A \lif !C$, kita kudu
nyathet instansi tautologis kang dawa iki:
\[
(!A \lif !B) \lif ((!B \lif !C) \lif (!A \lif !C))
\]
banjur ngetrapake \MP{} kaping pindho. Kita uga kerep pengin ngganti
sub-!!{formula} nganggo formula liyane kang wis dingerteni ekuivalen,
umpamane \iftag{prvDiamond}{$\Diamond
  !A$ nganggo $\lnot\Box\lnot !A$}{$\Box !A$ nganggo $\lnot\Diamond\lnot !A$},
utawa $\lnot\lnot !A$ nganggo $!A$. Mula tinimbang nulis kabeh
!!{derivation} nyata, luwih trep yen kita mung nyathet sebab apa
langkah antara iku !!{derivable}. Kanggo iku, ayo nglumpukake
sawetara fakta babagan !!{derivability}.

\begin{prop}
  Yen $\Log{K} \Proves !A_1$, \dots, $\Log{K} \Proves !A_n$, lan $!B$
  ngetutake saka $!A_1$, \dots, $!A_n$ miturut logika proposisional,
  mula $\Log{K} \Proves !B$.
\end{prop}

\begin{proof}
  Yen $!B$ ngetutake saka $!A_1$, \dots,~$!A_n$ miturut logika
  proposisional, mula
  \[
  !A_1 \lif (!A_2 \lif \cdots (!A_n \lif !B)\dots)
  \]
  iku instansi tautologis. Kanthi ngetrapake \MP{} kaping $n$, kita
  oleh !!{derivation} saka~$!B$.
\end{proof}

Kita bakal nandhani panganggone proposisi iki nganggo~\PL.

\begin{prop}
  If $\Log{K} \Proves !A_1 \lif (!A_2 \lif \cdots (!A_{n-1} \lif
  !A_n)\dots)$ then $\Log{K} \Proves \Box !A_1 \lif (\Box !A_2 \lif
  \cdots (\Box !A_{n-1} \lif \Box !A_n)\dots)$.
\end{prop}

\begin{proof}
  Kanthi induksi marang $n$, padha kaya bukti \olref[nor]{prop:rk}.
\end{proof}

Kita bakal nandhani panganggone proposisi iki nganggo~\RK.
Ayo dideleng kepiye asil iki nggampangake netepake
!!{derivability}.

\begin{prop}
  $\Log{K} \Proves (\Box!A \land \Box!B) \lif \Box (!A \land !B)$
\end{prop}

\begin{proof}
  \begin{derivation}
    1. & $\Log{K} \Proves !A \lif (!B \lif (!A \land !B))$ & \Taut \\
    2. & $\Log{K} \Proves \Box!A \lif (\Box !B \lif \Box(!A \land !B)))$ & \RK, 1\\
    3. & $\Log{K} \Proves (\Box!A \land \Box!B) \lif \Box(!A \land !B)$ & \PL, 2
  \end{derivation}
\end{proof}

\begin{prop}\ollabel{prop:rewriting}
  Yen $\Log{K} \Proves !A \liff !B$ lan $\Log{K} \Proves
  \Subst{!C}{!A}{q}$, mula $\Log{K} \Proves \Subst{!C}{!B}{q}$.
\end{prop}

\begin{proof}
  Gladhen.
\end{proof}

\begin{prob}
  Buktekna \olref[nml][prf][der]{prop:rewriting} kanthi induksi
  marang kerumitan~$!C$: yen $\Log{K} \Proves !A \liff !B$,
  mula $\Log{K} \Proves \Subst{!C}{!A}{q} \liff \Subst{!C}{!B}{q}$.
\end{prob}

Proposisi iki migunani banget nalika kita pengin ngowahi
$\Diamond$ dadi $\Box$ (utawa kosok baline), utawa mbusak negasi
dobel ing njero !!a{formula}. Ing ngisor iki, kita nandhani panganggone
\olref{prop:rewriting} nganggo ``$!B$ kanggo $!A$'' nalika nggunakake
!!a{formula}~$!C(!B)$ minangka panggantos kanggo~$!C(!A)$. Kanthi tembung liya,
``$!B$ kanggo $!A$'' iku cekakan saka:
  \begin{derivation}
    & $\Proves !C(!A)$ \\
    & $\Proves !A \liff !B$\\
    & $\Proves !C(!B)$ & miturut \olref{prop:rewriting}
  \end{derivation}
Upamane:

\begin{prop}
  $\Log{K} \Proves \lnot\Box p \lif \Diamond \lnot p$
\end{prop}

\begin{proof}
  \iftag{prvDiamond}{% Diamond iku primitif
  \begin{derivation}
    1. & $\Log{K} \Proves \Diamond \lnot p \liff \lnot\Box\lnot\lnot p$ &
    \Dual\\
    2. & $\Log{K} \Proves \lnot \Box \lnot\lnot p \lif \Diamond \lnot p$ & \PL, 1\\
    3. & $\Log{K} \Proves \lnot\Box p \lif \Diamond \lnot p$ & $p$ for $\lnot\lnot p$\\
  \end{derivation}
  }{% Diamond didéfinisikake
  \begin{derivation}
    1. & $\Log{K} \Proves \lnot \Box \lnot\lnot p \lif \Diamond \lnot p$ & \Taut\\
    2. & $\Log{K} \Proves \lnot\Box p \lif \Diamond \lnot p$ & $p$ for $\lnot\lnot p$\\
  \end{derivation}
  Formula ing baris~$1$ iku instansi saka $\lnot q \lif \lnot q$,
  kanthi substitusi $\Box\lnot\lnot p$ kanggo~$q$. $\Diamond\lnot p$ lan
  $\lnot\Box\lnot\lnot p$ padha, amarga $\Diamond$ didéfinisikake
  minangka~$\lnot\Box\lnot$.}
\end{proof}

Ing !!{derivation} ndhuwur, langkah pungkasan
``$p$ kanggo $\lnot\lnot p$'' iku cekakan saka
  \begin{derivation}
    & $\Log{K} \Proves \lnot \Box \lnot\lnot p \lif \Diamond \lnot
    p$\\
    & $\Log{K} \Proves \lnot\lnot p \liff p$ & \Taut\\
    & $\Log{K} \Proves \lnot\Box p \lif \Diamond \lnot p$ & miturut \olref{prop:rewriting}
  \end{derivation}
Ing kene $!C(q)$, $!A$, lan $!B$ saka \olref{prop:rewriting}
dipadhaake, miturut urutane, karo $\lnot \Box q \lif \Diamond \lnot p$,
$\lnot\lnot p$, lan~$p$.

Yen !!a{formula} ngemot sub-!!{formula}
$\lnot\Diamond !A$, kita bisa ngganti nganggo $\Box\lnot !A$
kanthi \olref{prop:rewriting}, amarga
$\Log{K} \Proves \lnot\Diamond !A \liff \Box\lnot !A$. Kita bakal
nandhani panggantos iki lan kang padha mung nganggo
``$\Box\lnot$ kanggo $\lnot\Diamond$.''

Proposisi iki mbenerake panganggone skema kanggo
netepake asil !!{derivability}. Upamane, proposisi sadurunge ora mung
netepake $\Log{K} \Proves \lnot\Box p \lif \Diamond \lnot p$,
nanging uga $\Log{K} \Proves \lnot\Box !A \lif \Diamond
\lnot !A$ kanggo sembarang~$!A$.

\begin{prop}
  Yen $!A$ instansi substitusi saka $!B$ lan $\Log{K} \Proves !B$,
  mula $\Log{K} \Proves !A$.
\end{prop}

\begin{proof}
  Mriksa manawa substitusi ing sakabehe !!{derivation} ngasilake
  !!{derivation} kang bener iku mbulet nanging lumrah, kanthi induksi
  marang dawa !!{derivation} saka~$!B$. Mligine, instansi substitusi
  saka instansi tautologis tetep instansi tautologis; instansi
  substitusi saka instansi~\iftag{prvDiamond}{\Dual{} lan~}{}\Ax{K}
  tetep instansi saka~\iftag{prvDiamond}{\Dual{}
    lan~}{}\Ax{K}; lan panganggone \MP{} lan~\Nec{} tetep bener
  nalika !!{formula}s disubstitusikake kanggo
  !!{propositional variable}s ing premis lan kesimpulane.
\end{proof}

\end{document}
